
\begin{abstract}
    Information about pretraining corpora used to train the current best-performing language models is seldom discussed:
    commercial models rarely detail their data, and
    even open models are often released without accompanying training data or recipes to reproduce them. 
    As a result, it is challenging to conduct and advance scientific research on language modeling, such as understanding how training data impacts model capabilities and limitations.
    To facilitate scientific research on language model pretraining, we curate and release \textbf{\dolma}, a three-trillion-token English corpus, built from a diverse mixture of web content, scientific papers, code, public-domain books, social media, and encyclopedic materials. 
    We extensively document
    \dolma, including its design principles, details about its construction, and a summary of its contents.
    We present analyses and experimental results on intermediate states of \dolma to share what we have learned about important data curation practices.
    Finally, we open-source our data curation toolkit
    to enable reproduction of our work as well as support further research in large-scale data curation.\footnotemark[1]

    \renewcommand{\arraystretch}{1.2}
    \begin{tabular}{rl}
         \huggingface & \href{https://huggingface.co/datasets/allenai/dolma}{\path{hf.co/datasets/allenai/dolma}}\\
         \github & \href{https://github.com/allenai/dolma}{\path{github.com/allenai/dolma}} \\
    \end{tabular}
    
    

    

    
\end{abstract}
